% Web専用草稿．独立査読未了．
% 文献調査：所蔵文献を検索したが該当箇所の本文は未照合．既出知識から独自に論証する．
\section{有限集合の数え上げ}
% genron:begin K-001204
% genron:use K-000001
% genron:use K-000002
% genron:use K-000003
% genron:use K-000004
% genron:use K-000005
% genron:use K-000007
% genron:use K-000017
% genron:use K-000024
% genron:use K-000110
% genron:use K-000115
% genron:use K-000118
\begin{prop}[加法原理]
互いに素な有限集合 $A_1,\ldots,A_k$ に対し $\left|\bigcup_{i=1}^kA_i\right|=\sum_{i=1}^k|A_i|$ である．
\end{prop}
\begin{proof}
各集合の元を一列に並べ，それらを順に連結すれば和集合と長さ $\sum_i|A_i|$ の自然数区間との全単射を得る．空の族では両辺は零である．
\end{proof}
% genron:end K-001204
% genron:begin K-001205
% genron:use K-000001
% genron:use K-000002
% genron:use K-000003
% genron:use K-000004
% genron:use K-000005
% genron:use K-000007
% genron:use K-000017
% genron:use K-000024
% genron:use K-001204
% genron:use K-000116
% genron:use K-000100
\begin{prop}[乗法原理]
有限集合 $A_1,\ldots,A_k$ に対し $\left|\prod_{i=1}^kA_i\right|=\prod_{i=1}^k|A_i|$ である．
\end{prop}
\begin{proof}
二集合の場合，直積を第一成分の値で互いに素な部分集合に分割すると各部分集合は第二集合と全単射である．加法原理を適用し，有限個の場合を帰納法で得る．空直積は一元集合である．
\end{proof}
% genron:end K-001205
% genron:begin K-001206
% genron:use K-000001
% genron:use K-000002
% genron:use K-000003
% genron:use K-000004
% genron:use K-000005
% genron:use K-000007
% genron:use K-000017
% genron:use K-000024
% genron:use K-000105
% genron:use K-000110
% genron:use K-000117
\begin{defn}[階乗と二項係数]
$0!=1$，$(n+1)!=(n+1)n!$ と定める．$n$ 元集合の $k$ 元部分集合の個数を $\binom nk$ と書き，$k>n$ では零とする．
\end{defn}
% genron:end K-001206
% genron:begin K-001207
% genron:use K-000001
% genron:use K-000002
% genron:use K-000003
% genron:use K-000004
% genron:use K-000005
% genron:use K-000007
% genron:use K-000017
% genron:use K-000024
% genron:use K-001204
% genron:use K-001206
% genron:use K-000112
% genron:use K-000118
% genron:use K-000100
\begin{prop}[順列の個数]
$n$ 元集合から異なる $k$ 個を順序付きで選ぶ方法の数は，$0\leq k\leq n$ のとき $n(n-1)\cdots(n-k+1)$ である．特に $n$ 元集合の全順列の個数は $n!$ である．
\end{prop}
\begin{proof}
長さ $k$ の列を最初の元で分割する．各分枝では残る $n-1$ 元から長さ $k-1$ の列を選ぶ．長さ零の列が一つであることを初期値とし，加法原理と帰納法によって積の式を得る．
\end{proof}
% genron:end K-001207
% genron:begin K-001208
% genron:use K-000001
% genron:use K-000002
% genron:use K-000003
% genron:use K-000004
% genron:use K-000005
% genron:use K-000007
% genron:use K-000017
% genron:use K-000024
% genron:use K-001207
% genron:use K-001206
% genron:use K-001204
\begin{prop}[二項係数の積表示]
$0\leq k\leq n$ に対して
\[
k!\binom nk=n(n-1)\cdots(n-k+1),\qquad
k!(n-k)!\binom nk=n!
\]
である．
\end{prop}
\begin{proof}
順序付きの $k$ 元列を，その列に現れる元の集合によって分割する．各 $k$ 元部分集合につき順列が $k!$ 通りある．順列の個数と加法原理により第一式を得る．両辺に $(n-k)!$ を掛ければ第二式となる．
\end{proof}
% genron:end K-001208
% genron:begin K-001209
% genron:use K-000001
% genron:use K-000002
% genron:use K-000003
% genron:use K-000004
% genron:use K-000005
% genron:use K-000007
% genron:use K-000017
% genron:use K-000024
% genron:use K-001206
% genron:use K-001204
\begin{prop}[Pascalの関係式]
$n\geq1$，$1\leq k\leq n$ について $\binom nk=\binom{n-1}{k}+\binom{n-1}{k-1}$ である．また $\binom n0=\binom nn=1$ である．
\end{prop}
\begin{proof}
一つの元 $a$ を固定して $k$ 元部分集合を $a$ を含まないものと含むものに分ける．前者は残りの $n-1$ 元集合の $k$ 元部分集合，後者はそこから $a$ を取り除けば $k-1$ 元部分集合と全単射になる．境界値は空集合と全体集合の一意性による．
\end{proof}
% genron:end K-001209
% genron:begin K-001210
% genron:use K-000001
% genron:use K-000002
% genron:use K-000003
% genron:use K-000004
% genron:use K-000005
% genron:use K-000007
% genron:use K-000017
% genron:use K-000024
% genron:use K-001209
% genron:use K-000249
% genron:use K-000100
\begin{prop}[二項定理]
可換環の元 $a,b$ と自然数 $n$ に対して $(a+b)^n=\sum_{k=0}^n\binom nk a^{n-k}b^k$ である．
\end{prop}
\begin{proof}
$n=0$ は単位元の定義による．$n$ での式に $a+b$ を掛け，分配律と可換性を用いて同じ冪の項をまとめると内側の係数は Pascal の関係式，端の係数は一となる．数学的帰納法で成立する．
\end{proof}
% genron:end K-001210
